
\documentclass[11pt,a4paper]{article}

%------------------------------------------------
% Encoding and language
%------------------------------------------------
\usepackage[T1]{fontenc}
\usepackage[utf8]{inputenc}
\usepackage[english]{babel}

%------------------------------------------------
% Page layout
%------------------------------------------------
\usepackage[a4paper,margin=2.8cm]{geometry}
\usepackage{microtype}
\usepackage{setspace}
\setstretch{1.1}

%------------------------------------------------
% Mathematics
%------------------------------------------------
\usepackage{amsmath,amssymb,amsfonts,amsthm}
\usepackage{mathtools}
\usepackage{mathrsfs}
\usepackage{bm}
\usepackage{enumitem}

%------------------------------------------------
% Hyperlinks
%------------------------------------------------
\usepackage[colorlinks=true,
linkcolor=blue,
citecolor=blue,
urlcolor=blue]{hyperref}

%------------------------------------------------
% Theorem environments
%------------------------------------------------
\theoremstyle{plain}
\newtheorem{theorem}{Theorem}[section]
\newtheorem{conjecture}[theorem]{Conjecture}
\newtheorem{proposition}[theorem]{Proposition}
\newtheorem{lemma}[theorem]{Lemma}
\newtheorem{corollary}[theorem]{Corollary}

\theoremstyle{definition}
\newtheorem{definition}[theorem]{Definition}
\newtheorem{example}[theorem]{Example}

\theoremstyle{remark}
\newtheorem{remark}[theorem]{Remark}

%------------------------------------------------
% Mathematical notation
%------------------------------------------------

% Number systems
\newcommand{\R}{\mathbb{R}}
\newcommand{\C}{\mathbb{C}}
\newcommand{\Z}{\mathbb{Z}}
\newcommand{\N}{\mathbb{N}}
\newcommand{\Q}{\mathbb{Q}}

% Symplectic geometry
\newcommand{\M}{M}
\newcommand{\Ham}{\operatorname{Ham}}
\newcommand{\Symp}{\operatorname{Symp}}
\newcommand{\Diff}{\operatorname{Diff}}
\newcommand{\id}{\operatorname{id}}
\newcommand{\Fix}{\operatorname{Fix}}
\newcommand{\Crit}{\operatorname{Crit}}
\newcommand{\CritMorse}{\operatorname{Crit}_{\mathrm{Morse}}}
\newcommand{\supp}{\operatorname{supp}}

% Differential operators and sets
\newcommand{\d}{\mathrm{d}}
\newcommand{\Lie}{\mathcal{L}}
\newcommand{\grad}{\nabla}
\newcommand{\eps}{\varepsilon}
\newcommand{\Om}{\Omega}

%------------------------------------------------
% Title information
%------------------------------------------------
\title{\textbf{Arnold's Conjecture on Fixed Points of
		Hamiltonian Diffeomorphisms}}

\author{}
\date{}

%================================================
% Document
%================================================
\begin{document}
	
	\maketitle
	
	%------------------------------------------------
	% Abstract
	%------------------------------------------------
	\begin{abstract}
		Arnold's conjecture is a fundamental problem in symplectic
		geometry that relates the fixed points of Hamiltonian
		diffeomorphisms to the topology of the underlying symplectic
		manifold. More precisely, it predicts that the number of
		fixed points of a Hamiltonian diffeomorphism is bounded below
		by the minimal number of critical points of a smooth function
		on the manifold. When all fixed points are nondegenerate,
		the conjecture asserts the stronger lower bound given by the
		minimal number of critical points of a Morse function.
		
		The objective of this work is to study Arnold's conjecture
		and the geometric and topological mechanisms underlying these
		lower bounds. We begin by introducing the basic notions of
		symplectic geometry and Hamiltonian dynamics, and then
		formulate the conjecture in terms of fixed points of
		Hamiltonian flows. We aim to understand how variational
		methods, Morse theory, and Floer homology can be used to
		relate the dynamics of Hamiltonian diffeomorphisms to the
		topology of the ambient manifold. Particular attention is
		given to the role of nondegeneracy and to the way in which
		topological invariants provide lower bounds for the number
		of fixed points.
	\end{abstract}
	
	\noindent
	\textbf{Keywords:}
	Arnold's conjecture, symplectic geometry, Hamiltonian
	diffeomorphisms, fixed points, Morse theory, Floer homology.
	
	\medskip
	
	\noindent
	\textbf{Mathematics Subject Classification (2020):}
	53D40, 37J10, 58E05.
	
	%------------------------------------------------
	\section{Introduction}
	%------------------------------------------------
	
	The study of fixed points of Hamiltonian diffeomorphisms is
	one of the central themes in symplectic geometry. Unlike
	general smooth dynamical systems, Hamiltonian systems possess
	a rich geometric structure determined by a symplectic form.
	This structure imposes strong restrictions on their dynamics
	and creates a natural connection between dynamical properties
	and the topology of the underlying manifold.
	
	Let $(M,\omega)$ be a closed symplectic manifold, where
	$\omega$ is a smooth, closed, nondegenerate differential
	$2$-form. A time-dependent Hamiltonian
	\[
	H\colon [0,1]\times M\longrightarrow \R
	\]
	generates a time-dependent vector field $X_{H_t}$ defined by
	the convention
	\[
	\iota_{X_{H_t}}\omega=-\d H_t,
	\qquad H_t(x):=H(t,x).
	\]
	The corresponding Hamiltonian flow is denoted by
	$\phi_H^t$, with $\phi_H^0=\id_M$. Its time-one map
	\[
	\phi_H:=\phi_H^1
	\]
	is a Hamiltonian diffeomorphism of $M$. We are interested
	in the fixed-point set
	\[
	\Fix(\phi_H)
	:=
	\{x\in M\mid \phi_H(x)=x\}.
	\]
	
	Arnold's conjecture, formulated by Vladimir Arnold, predicts
	that the number of fixed points of $\phi_H$ is controlled
	from below by topological invariants of $M$. In particular,
	the conjecture suggests that Hamiltonian dynamics cannot
	be arbitrarily simple: the topology of the manifold forces
	the existence of a certain number of fixed points.
	
	%------------------------------------------------
	\section{Statement of Arnold's Conjecture}
	%------------------------------------------------
	
	For a smooth function $f\colon M\to\R$, define its critical
	set by
	\[
	\Crit(f)
	:=
	\{x\in M\mid \d f(x)=0\}.
	\]
	We introduce the minimal number of critical points
	\[
	\Crit(M)
	:=
	\min_{f\in C^\infty(M,\R)}
	\#\Crit(f),
	\]
	where the minimum is taken over smooth functions with
	finitely many critical points.
	
	Similarly, let
	\[
	\CritMorse(M)
	:=
	\min\bigl\{
	\#\Crit(f)\mid
	f\in C^\infty(M,\R)
	\text{ is a Morse function}
	\bigr\}.
	\]
	
	The conjecture can then be stated as follows.
	
	\begin{conjecture}[Arnold]
		\label{conj:arnold}
		Let $(M,\omega)$ be a closed symplectic manifold, and let
		$\phi\colon M\to M$ be the time-one map of a
		time-dependent Hamiltonian flow. Then
		\[
		\boxed{
			\#\Fix(\phi)\geq \Crit(M).
		}
		\]
		If all fixed points of $\phi$ are nondegenerate, then
		\[
		\boxed{
			\#\Fix(\phi)\geq \CritMorse(M).
		}
		\]
	\end{conjecture}
	
	Here, a fixed point $x\in\Fix(\phi)$ is called
	\emph{nondegenerate} if
	\[
	\det\bigl(D\phi(x)-I_{T_xM}\bigr)\neq 0,
	\]
	where $D\phi(x)\colon T_xM\to T_xM$ is the derivative
	of $\phi$ at $x$ and $I_{T_xM}$ denotes the identity
	map on the tangent space.
	
	%------------------------------------------------
	\section{Objective and Strategy}
	%------------------------------------------------
	
	The principal objective of this work is to investigate
	Arnold's conjecture by exploring the relationship between
	Hamiltonian dynamics, variational methods, and the topology
	of closed symplectic manifolds.
	
	More specifically, we aim to address the following questions:
	
	\begin{enumerate}[label=\textnormal{(\arabic*)}]
		\item \textbf{Hamiltonian dynamics.}
		How can the fixed points of a Hamiltonian diffeomorphism
		be characterized in terms of periodic orbits of the
		associated Hamiltonian flow?
		
		\item \textbf{Morse theory.}
		Why does the topology of $M$ impose a lower bound on
		the number of critical points of a smooth function,
		and how does this motivate the corresponding lower
		bound for Hamiltonian fixed points?
		
		\item \textbf{Variational formulation.}
		Under which assumptions can periodic Hamiltonian
		orbits be interpreted as critical points of an action
		functional on a suitable loop space?
		
		\item \textbf{Floer homology.}
		How can an infinite-dimensional analogue of Morse
		theory be developed to establish lower bounds for
		the number of periodic orbits and fixed points?
		
		\item \textbf{Nondegeneracy.}
		What additional information becomes available when
		all fixed points are nondegenerate, and why is the
		Morse-theoretic lower bound relevant in this setting?
	\end{enumerate}
	
	Our approach is to develop the necessary background in
	symplectic geometry, review the relevant ideas from Morse
	theory, and then examine the role of Hamiltonian action
	functionals and Floer homology. The precise results that
	can be established depend on the hypotheses imposed on
	the symplectic manifold and on the Hamiltonian.
	
	%------------------------------------------------
	\section{Organization of the Work}
	%------------------------------------------------
	
	The work is organized as follows. Section~2 introduces
	the basic definitions and properties of symplectic
	manifolds and Hamiltonian flows. Section~3 reviews the
	Morse-theoretic background and explains the topological
	lower bounds appearing in Arnold's conjecture. Section~4
	develops the variational perspective and introduces the
	Hamiltonian action functional. Section~5 discusses the
	main ideas of Floer homology and their application to
	the fixed-point problem. Finally, we summarize the
	conclusions and indicate possible extensions.
	
	%------------------------------------------------
	% Bibliography
	%------------------------------------------------
	\begin{thebibliography}{99}
		
		\bibitem{Arnold}
		V.~I.~Arnold,
		\emph{Mathematical Methods of Classical Mechanics},
		2nd ed., Graduate Texts in Mathematics, vol.~60,
		Springer, New York, 1989.
		
		\bibitem{AudinDamian}
		M.~Audin and M.~Damian,
		\emph{Morse Theory and Floer Homology},
		Universitext, Springer, London, 2014.
		
		\bibitem{McDuffSalamon}
		D.~McDuff and D.~Salamon,
		\emph{Introduction to Symplectic Topology},
		3rd ed., Oxford University Press, Oxford, 2017.
		
		\bibitem{Salamon}
		D.~Salamon,
		\emph{Lectures on Floer Homology},
		in Symplectic Geometry and Topology,
		IAS/Park City Mathematics Series, vol.~7,
		American Mathematical Society, 1999.
		
	\end{thebibliography}
	
\end{document}
